\section{Implementacja przygotowania danych}

Przygotowano moduł odpowiedzialny za wczytanie danych, rozdzielenie cech wejściowych i etykiet, wykonanie podziału na zbiory uczące oraz standaryzację danych.

Zastosowane rozwiązanie zostało zaimplementowane w pliku \texttt{Code/scripts/dataloader.py}. Wydzielenie obsługi danych do osobnego modułu pozwala oddzielić etap przygotowania danych od implementacji samej sieci neuronowej. Dzięki temu kod modelu może korzystać z gotowych tablic cech i etykiet, bez konieczności bezpośredniego odczytywania pliku CSV.

\subsection{Wczytywanie zbioru danych}

Dane wejściowe są wczytywane z pliku CSV za pomocą funkcji \texttt{load\_wine\_dataset}. Funkcja przyjmuje ścieżkę do pliku z danymi oraz parametry określające proporcje podziału na zbiory treningowy, walidacyjny i testowy.

\begin{lstlisting}[language=Python, caption={Fragment funkcji wczytującej dane ze zbioru Wine.}, label={lst:load_wine_dataset_basic}, basicstyle=\ttfamily\small, breaklines=true, frame=single]
df = pd.read_csv(data_path, header=None)

# Pierwsza kolumna zawiera etykiety klas
# pozostale kolumny zawieraja cechy wejsciowe
y_raw = df.iloc[:, 0].to_numpy(dtype=np.int64)
x_raw = df.iloc[:, 1:].to_numpy(dtype=np.float64)
\end{lstlisting}

Takie rozwiązanie jest zgodne ze strukturą zbioru Wine.

\subsection{Podział na cechy wejściowe i etykiety}

Po wczytaniu pliku dane są dzielone na macierz cech wejściowych $X$ oraz wektor etykiet klas $y$. Macierz $X$ zawiera wartości wykorzystywane jako wejście sieci neuronowej, natomiast wektor $y$ zawiera informację o oczekiwanej klasie próbki. Rozdzielenie tych dwóch części jest konieczne, ponieważ podczas uczenia model otrzymuje na wejściu cechy próbki, a następnie porównuje uzyskaną odpowiedź z etykietą rzeczywistą.

\subsection{Kodowanie etykiet klas}

Oryginalne etykiety klas w zbiorze Wine są zapisane jako wartości $1$, $2$ i $3$. W implementacji zostały one przemapowane na zakres $0$, $1$, $2$. Takie kodowanie jest wygodniejsze, ponieważ indeks klasy może być bezpośrednio wykorzystany przy obliczaniu funkcji kosztu lub podczas tworzenia reprezentacji \textit{one-hot}

\begin{lstlisting}[language=Python, caption={Mapowanie etykiet klas na numerację od zera.}, label={lst:map_labels}, basicstyle=\ttfamily\small, breaklines=true, frame=single]
class_labels = sorted(np.unique(y).tolist())
mapping = {label: index for index, label in enumerate(class_labels)}
mapped = np.array([mapping[label] for label in y], dtype=np.int64)
\end{lstlisting}

Dla rozważanego zbioru to oznacza: klasa $1$ zostaje zapisana jako $0$, klasa $2$ jako $1$, a klasa $3$ jako $2$.

\subsection{Podział na zbiory treningowy, walidacyjny i testowy}

W module przygotowania danych zaimplementowano podział stratyfikowany, ponieważ liczebność klas w zbiorze Wine nie jest identyczna. Zastosowanie zwykłego losowego podziału mogłoby spowodować, że jedna z klas byłaby zbyt słabo reprezentowana w zbiorze testowym/walidacyjnym.

Podstawowy wariant funkcji umożliwia podział na zbiór treningowy i testowy. Dodatkowo przewidziano możliwość wydzielenia zbioru walidacyjnego przez ustawienie parametru \texttt{val\_size} > $0$. W testowym uruchomieniu zastosowano proporcje $70\% / 15\% / 15\%$, odpowiednio zbiór: treningowy, walidacyjny i testowy.

\begin{table}[htbp]
    \centering
    \caption{Rozmiary zbiorów uzyskane podczas testowego uruchomienia modułu DataLoader.}
    \label{tab:dataloader_split_shapes}
    \begin{tabular}{lcc}
        \toprule
        Zbiór & Kształt macierzy cech & Kształt wektora etykiet \\
        \midrule
        treningowy & $(125, 13)$ & $(125,)$ \\
        walidacyjny & $(26, 13)$ & $(26,)$ \\
        testowy & $(27, 13)$ & $(27,)$ \\
        \bottomrule
    \end{tabular}
\end{table}

\subsection{Standaryzacja danych}

Dla każdej cechy obliczana jest średnia $\mu$ oraz odchylenie standardowe $\sigma$. Następnie wartości są przekształcane zgodnie ze wzorem:

\begin{equation}
    x' = \frac{x - \mu}{\sigma}.
\end{equation}

Parametry standaryzacji są wyznaczane wyłącznie na podstawie zbioru treningowego. Następnie te same wartości średnich i odchyleń standardowych są stosowane do przeskalowania zbioru walidacyjnego i testowego. Zapobiega to przeciekowi informacji ze zbiorów wykorzystywanych do oceny modelu do etapu uczenia.

\begin{lstlisting}[language=Python, caption={Standaryzacja cech na podstawie zbioru treningowego.}, label={lst:standardization}, basicstyle=\ttfamily\small, breaklines=true, frame=single]
mean = np.mean(x_train, axis=0)
std = np.std(x_train, axis=0)
std[std == 0.0] = 1.0

x_train_std = (x_train - mean) / std
x_val_std = (x_val - mean) / std
x_test_std = (x_test - mean) / std
\end{lstlisting}

W implementacji dodano również zabezpieczenie przed dzieleniem przez zero.

\subsection{Generowanie mini-batchy}

Do obsługi danych podczas uczenia przygotowano klasę \texttt{BatchLoader}, której zadaniem jest dzielenie zbioru danych na mniejsze partie - mini-batche. Umożliwia to uczenie sieci na fragmentach danych zamiast przetwarzania całego zbióru treningowego jednocześnie.

Klasa \texttt{BatchLoader} przyjmuje macierz cech, wektor etykiet, rozmiar batcha oraz informację o tym, czy dane mają zostać przetasowane przed rozpoczęciem iteracji. Tasowanie danych jest przydatne w zbiorze treningowym, ponieważ ogranicza wpływ kolejności próbek na przebieg uczenia.

\begin{lstlisting}[language=Python, caption={Zwracanie kolejnych mini-batchy.}, label={lst:batch_loader}, basicstyle=\ttfamily\small, breaklines=true, frame=single]
for start in range(0, len(indices), self.batch_size):
    batch_indices = indices[start:start + self.batch_size]
    yield self.x[batch_indices], self.y[batch_indices]
\end{lstlisting}

Podczas testowego uruchomienia dla parametru \texttt{batch\_size = 16} pierwszy batch treningowy miał kształt $(16, 13)$ dla macierzy cech oraz $(16,)$ dla wektora etykiet.

\subsection{Testy modułu przygotowania danych}

Poprawność działania modułu sprawdzono za pomocą skryptu \texttt{Code/test\_dataloader.py}. Skrypt wczytuje ścieżkę do zbioru danych ze zmiennej środowiskowej \texttt{WINE\_DATA\_PATH}, uruchamia funkcję \texttt{load\_wine\_dataset}, a następnie wypisuje podstawowe informacje o otrzymanych zbiorach. Dodatkowo sprawdzane jest działanie klasy \texttt{BatchLoader} przez wypisanie kształtu pierwszego batcha treningowego i walidacyjnego.

